% Part: applied-modal-logics
% Chapter: temporal-logic
% Section: possible-histories

\documentclass[../../../include/open-logic-section]{subfiles}

\begin{document}

\olfileid[pt]{aml}{tl}{poss}

\olsection{Histórias possíveis}

Os modelos relacionais de lógica temporal que temos usado são
extremamente flexíveis, pois não precisamos impor restrição alguma à
relação de acessibilidade. Isso significa que os modelos temporais
podem ramificar-se no passado e no futuro, mas talvez queiramos
considerar uma concepção mais ``modal'' de ramificação, na qual
sequências de acontecimentos são tratadas como histórias possíveis.
Isso não exige necessariamente uma mudança de linguagem, embora
também possamos acrescentar os operadores modais ``usuais'' $\Box$
e~$\Diamond$, bem como relações de acessibilidade epistêmica para
representar mudanças no conhecimento dos agentes ao longo do tempo.

\begin{defn}
  Um \emph{modelo de histórias possíveis} para a linguagem temporal é uma tripla
  $\mModel{M} = \tuple{T, C, V}$, em que
  \begin{enumerate}
    \item $T$ é um conjunto não vazio, interpretado como estados temporais.
    \item $C$ é um conjunto de trajetórias computacionais, ou
      \emph{histórias possíveis}, de um sistema. Em outras palavras,
      $C$~é um conjunto de sequências~$\sigma$ de estados $s_1$,
      $s_2$,~$s_3$, \dots, em que todo $s_i \in T$.
    \item $V$ é uma função que atribui a cada !!{propositional
      variable}~$p$ um conjunto~$V(p)$ de instantes temporais.
  \end{enumerate}
  Para simplificar, em geral também suporemos que, quando uma história
  pertence a~$C$, todos os seus sufixos também pertencem. Por exemplo,
  se $s_1$, $s_2$,~$s_3$ é uma sequência em~$C$, então $s_2$, $s_3$
  e~$s_3$ também o são. Além disso, quando dois estados $s_i$ e~$s_j$
  aparecem em uma sequência~$\sigma$, dizemos que
  $s_i \prec_\sigma s_j$ quando $i < j$. Quando $t \in V(p)$, dizemos
  que $p$~é \emph{verdadeira em}~$t$.
\end{defn}

A única mudança relevante é que, ao avaliarmos a verdade de
!!a{formula} em um instante~$t$ de um modelo~$\mModel M$, fazemos isso
em relação a uma história~$\sigma$, na qual $t$ aparece como estado.
Não precisamos, contudo, mudar a semântica das !!{propositional
variable}s nem a dos conectivos vero-funcionais. Elas permanecem
exatamente como em \olref[sem]{defn:tmodels}, pois nenhuma faz
referência a~$\sigma$. Agora, porém, redefinimos o operador de
futuro~$\Ftemp$ e acrescentamos o operador~$\Diamond$ em relação a
essas histórias. 

\begin{defn}\ollabel{defn:phmodels}
  A \emph{verdade de !!a{formula}~$!A$ em~$t, \sigma$} em~$\mModel M$, em símbolos:
  $\mSat{M}{!A}[t, \sigma]$:
  \begin{enumerate}
  \item\ollabel{defn:sub:phmodels-f} \indcase{!A}{\Ftemp
    !B}{$\mSat{M}{\indfrm}[t, \sigma]$ sse $\mSat{M}{!B}[t', \sigma]$
    para algum $t' \in T$ tal que $t \prec_\sigma t'$.}
  \item\ollabel{defn:sub:phmodels-diamond} \indcase{!A}{\Diamond
    !B}{$\mSat{M}{\indfrm}[t, \sigma]$ sse $\mSat{M}{!B}[t, \sigma']$
    para algum $\sigma' \in C$ no qual $t$ ocorra.}
  \end{enumerate} 
\end{defn}

Outros operadores temporais e modais podem ser definidos de maneira
semelhante. Agora, porém, podemos representar afirmações que combinam
tempo verbal e modalidade. Por exemplo, podemos simbolizar
``$p$~não ocorrerá, mas poderia ter ocorrido'' pela fórmula
$\lnot \Ftemp p \land \Diamond \Ftemp p$. Ela valeria em um instante e
uma história nos quais $p$ não se torna verdadeira em um estado
sucessor, mas existe uma história alternativa na qual $p$ se tornará
verdadeira. 

\end{document}
